\chapter{What was added or corrected relative to the shorter theory note}

\begin{enumerate}[leftmargin=1.6em]
\item The local Taylor expansion, selection gradient, whitening transformation, symmetric operator, generalized eigenproblem, and mode solution are derived line by line.
\item The fixed-\(\G\) assumption behind \(\dot{\mathbf x}=-\Sop\mathbf x\) is stated explicitly; the moving-metric term is identified when \(\G(t)\) changes.
\item Channel comparison for \(\G_1\neq\G_2\) is proven coordinate invariant under a common metric, while the scientific choice of that metric is distinguished from mere coordinate invariance.
\item The projector proof of the fixed-representation local quadratic law is supplied in full.
\item A second theorem derives the fully invariant coefficient \(\mathbf K_M\) when \(\G\) evolves, including both rotation of the \(\Sop\) eigenspace and direct deformation by \(\G^{1/2}\).
\item The role of commutation and off-diagonal forcing is separated from the role of the spectral gap; a diagonal example shows explicitly that a positive Davis--Kahan bound can accompany zero actual rotation.
\item The deterministic-versus-stochastic result is separated into an exact tangent-surrogate identity \(a t+b t^2\) and the actual chordal asymptotic \(a t+O(t^2)\), whose second-order term can contain stochastic curvature corrections.
\item Expected log fitness and arithmetic mean fitness are distinguished and each is derived separately.
\item The general trace identity and channel-ridge theorem are proven in full, including the condition required for the ridge to be the leading response channel.
\item The standing-variation derivation now includes the migration contribution
\[
\dot{\zbar}_i=\G_i\betav_i+m(\zbar_j-\zbar_i).
\]
\item The derivative of \(\G^{1/2}\) is solved explicitly in the eigenbasis of \(\G\), allowing \(\dot\Sop\) to be partitioned into genetic and curvature components.
\item The DMI handoff is given a Poisson interpretation and its early-time expansion clarifies what it does and does not change about the leading exponent.
\item The existing covariance-tensor machinery is derived through the isometric \(\svec\) map and connected explicitly to the new response-operator geometry.
\end{enumerate}

\clearpage
\begin{thebibliography}{99}

\bibitem{aguirre2014}
Aguirre, J. D., Hine, E., McGuigan, K. and Blows, M. W. (2014).
Comparing \(\G\): multivariate analysis of genetic variation in multiple populations.
\emph{Heredity} \textbf{112}, 21--29.

\bibitem{bjorckgolub1973}
Bj\"orck, \AA. and Golub, G. H. (1973).
Numerical methods for computing angles between linear subspaces.
\emph{Mathematics of Computation} \textbf{27}, 579--594.

\bibitem{daviskahan1970}
Davis, C. and Kahan, W. M. (1970).
The rotation of eigenvectors by a perturbation. III.
\emph{SIAM Journal on Numerical Analysis} \textbf{7}, 1--46.

\bibitem{hine2009}
Hine, E., Chenoweth, S. F., Rundle, H. D. and Blows, M. W. (2009).
Characterizing the evolution of genetic variance using genetic covariance tensors.
\emph{Philosophical Transactions of the Royal Society B} \textbf{364}, 1567--1578.

\bibitem{lande1979}
Lande, R. (1979).
Quantitative genetic analysis of multivariate evolution, applied to brain:body size allometry.
\emph{Evolution} \textbf{33}, 402--416.

\bibitem{landearnold1983}
Lande, R. and Arnold, S. J. (1983).
The measurement of selection on correlated characters.
\emph{Evolution} \textbf{37}, 1210--1226.

\bibitem{lynchwalsh1998}
Lynch, M. and Walsh, B. (1998).
\emph{Genetics and Analysis of Quantitative Traits}.
Sinauer Associates, Sunderland, Massachusetts.

\bibitem{orr1995}
Orr, H. A. (1995).
The population genetics of speciation: the evolution of hybrid incompatibilities.
\emph{Genetics} \textbf{139}, 1805--1813.

\bibitem{schluter1996}
Schluter, D. (1996).
Adaptive radiation along genetic lines of least resistance.
\emph{Evolution} \textbf{50}, 1766--1774.

\bibitem{walshlynch2018}
Walsh, B. and Lynch, M. (2018).
\emph{Evolution and Selection of Quantitative Traits}.
Oxford University Press, Oxford.

\bibitem{walter2016}
Walter, G. M., Wilkinson, M. J., James, M. E., Richards, T. J., Aguirre, J. D. and Ortiz-Barrientos, D. (2016).
Diversification across a heterogeneous landscape.
\emph{Evolution} \textbf{70}, 1979--1992.

\bibitem{walter2018}
Walter, G. M., Aguirre, J. D., Blows, M. W. and Ortiz-Barrientos, D. (2018).
Evolution of genetic variance during adaptive radiation.
\emph{The American Naturalist} \textbf{191}, E108--E128.

\bibitem{walter2020}
Walter, G. M., Richards, T. J., Wilkinson, M. J., Blows, M. W., Aguirre, J. D. and Ortiz-Barrientos, D. (2020).
Loss of ecologically important genetic variation in late generation hybrids reveals links between adaptation and speciation.
\emph{Evolution Letters} \textbf{4}, 302--316.

\end{thebibliography}
